\section{Adversarial Pressure、Bias 与 Non-Authority}

真实用户不会像 benchmark annotator 那样配合任务定义。公开 demo 一上线，用户会寻找边界、制造两害比较、替换身份词并故意组合罕见情境。本章把这些攻击视为 evidence，而不是 PR 插曲。

\subsection{Exposed Socket：正确例子也要解释为什么}

新闻页描述 Alexa 让儿童触碰插座硬币挑战，右侧 Delphi 判断“让孩子碰裸露插座是危险的”。这个例子展示系统可以连接自然语言与物理后果，但它只是一个成功案例。真正的 evaluation 需要改写动作主体、时间、否定和 hypothetical context，并检查模型是否仍保持 danger concept。

\lecturefigure{slide-36-alexa-socket-delphi.jpg}{Alexa 暴露插座事件与 Delphi 的危险判断}{官方录像恢复幻灯片 V036；Stanford CS25 Lecture 16}

\subsection{上线几小时后，任务定义被用户重写}

backlash slide 记录 10 月 15 日公开后的变化：relative comparison mode 在数小时内被移除，一个周末收到约 25,000 个 adversarial examples，随后野生查询达到数百万。读图时应把这些数字理解为 attack-surface evidence；它们不等于均匀测试集，却揭示了 benchmark 很难预先枚举的组合空间。

\lecturefigure{slide-37-delphi-backlash-scale.jpg}{Delphi 上线后的快速下线、25k adversarial examples 与野生流量}{官方录像恢复幻灯片 V037；Stanford CS25 Lecture 16}

下一页展示研究者和公众故意构造反常句子来“折磨”模型。此类例子常缺少现实语境，却非常擅长暴露 lexical shortcut。评价时不应在“题目荒谬所以失败不算”和“一个失败所以系统完全无用”之间二选一，而应建立 failure taxonomy：negation、role reversal、rare composition、identity substitution 与 nonsensical premise。

\lecturefigure{slide-38-adversarial-twitter-example.jpg}{社交媒体上的 contrived adversarial prompt}{官方录像恢复幻灯片 V038；Stanford CS25 Lecture 16}

\subsection{Public Critique 不是可以忽略的噪声}

《纽约时报》“Can a Machine Learn Morality?”与龙形插画把争议推向公共领域。下一页列出 concerned voices，并回到论文题目。两页共同说明 machine ethics 的评价者不只包括 ML reviewer，还包括伦理学者、记者、被模型影响的群体和普通用户；系统文档必须让这些人看见训练目标与限制。

\lecturefigure{slide-39-can-machine-learn-morality-nyt.jpg}{媒体追问：机器能否学习道德}{官方录像恢复幻灯片 V039；Stanford CS25 Lecture 16}

\lecturefigure{slide-40-concerned-voices.jpg}{对 Delphi 的公共担忧与研究反思}{官方录像恢复幻灯片 V040；Stanford CS25 Lecture 16}

\subsection{Delphi 不是 Moral Authority}

non-authority slide 明确写出：研究团队不支持用 AI 给出 moral advice，Delphi 不会替代法庭中的 human judges；下方引用说明项目目标是理解 imperfect humans 如何作伦理判断，而不是让机器成为终极道德权威。这个边界应进入任何下游使用许可和 UI，而不是只留在论文 discussion。

\lecturefigure{slide-41-no-ai-moral-authority.jpg}{明确边界：Delphi 不提供 moral advice，也不替代人类法官}{官方录像恢复幻灯片 V041；Stanford CS25 Lecture 16}

\begin{importantbox}{Capability Claim 与 Authority Claim 必须分开}
模型能预测一部分人类判断，是 capability claim；模型应据此决定惩罚、医疗、招聘或司法结果，是 authority claim。前者即使成立，也不能自动推出后者。权限需要制度、责任、申诉和审计设计。
\end{importantbox}

\subsection{“不做研究”也有 Normative Consequence}

下一页先承认 moral models 太困难、准确率有限，随后指出当前 AI 系统已经在做隐含伦理决定。最后的 callout 主张投资研究并形成 human-centered、unbiased、robust 的社会规范。读图时应看两侧风险：部署 imperfect moral model 会伤害人，完全不研究也可能把未经审计的默认规则留在系统中。

\lecturefigure{slide-42-ethics-investment-needed.jpg}{困难不等于可以忽略：现有系统已在做带伦理后果的决定}{官方录像恢复幻灯片 V042；Stanford CS25 Lecture 16}

“Delphi will empower powerful people and harm marginalized people”是反方担忧；后续 build 指出现有系统已经可能 reinforce unjust status quo。关键不是哪一条气泡胜出，而是部署评估必须比较 intervention 与 baseline：新系统造成的增量伤害是什么，现有流程又在伤害谁？

\lecturefigure{slide-43-reinforce-status-quo.jpg}{Intervention risk 与 status-quo risk 的双向比较}{官方录像恢复幻灯片 V043；Stanford CS25 Lecture 16}

\subsection{Bias Acknowledgement 要具体到 Data 与 Use}

bias slide 明确写出 norm bank 反映 2020--2021 annotation data 的 ideologies，judgments 由美国 crowdworkers 提供并偏向 Western-centric viewpoints，定制 Delphi 到其他文化或 moral frameworks 是开放方向。左侧新闻标题“Delphi leans left”提醒政治倾向也会被测量与传播。

\lecturefigure{slide-44-delphi-bias-acknowledgement.jpg}{Delphi 的数据、政治与文化偏差声明}{官方录像恢复幻灯片 V044；Stanford CS25 Lecture 16}

\teachervoice{讲者没有把 bias acknowledgement 当成免罪声明。她说 Delphi 的确 left-leaning、Western-centric，随后把问题转向更广泛参与和数据共享。承认偏差是审计起点，不是部署许可。}

\subsection{本章小结}

公开 demo 把 Delphi 从静态 benchmark 推入 adversarial、政治与制度环境。最重要的课堂边界是 non-authority：预测描述性判断不授予决策权。可靠评估必须同时记录模型失败、用户攻击、数据群体与 status-quo baseline。

